\section{Introduction}
\label{sec:intro}

An attacker who has just landed inside a honeypot issues a single
request: \emph{show me the database credentials}. What the honeypot does
next divides two decades of deception engineering from the problem this
paper addresses. A recorder cannot answer at all; it writes the request
to disk and waits. A scripted honeypot answers from a file, the same file
for every attacker, on every visit, forever. A generative honeypot drafts
an answer on the spot---and the attacker may act on it, carrying it into
a timeline a security operations center will later treat as evidence. A
drafted credential exists nowhere except in the attacker's hands, and
nothing downstream can tell it apart from a secret that leaked; an
echoed real one converts the trap into a breach. The moment a honeypot
begins to answer it has taken on three jobs at once: stay consistent
with everything it has said, leak nothing real, and remain available
when the model behind it is not. A recent survey of production deception frameworks
identifies a gap in the surveyed systems: responses are not grounded in
real-time threat intelligence and behavior is not differentiated across
attacker classes~\cite{lopez2024}.

The cost of that silence is measured elsewhere: mean time to identify a
breach runs in months~\cite{ibm2024} while adversary breakout time has
compressed to minutes~\cite{crowdstrike2026}. A honeypot sits directly in
the attacker's path yet converts none of that proximity into
tactics-and-techniques intelligence a security operations center could act
on.

Large language models (LLMs) promise to close this gap. Prompted with a
session transcript and a library of threat intelligence, a modern LLM can
impersonate a vulnerable Linux host, improvise plausible file systems, and
steer the conversation toward the attacker's next objective. Recent systems
have begun to exploit this, deploying LLM agents as honeypots that chat an
attacker through a fictitious machine~\cite{reworr2024}. The promise,
however, arrives with a problem the deception literature has largely left
unexamined: \emph{the model can be wrong in ways a script cannot}. A
honeypot's product is a lie, and the lie must satisfy three properties at
once. It must be \emph{consistent}: the kernel version printed on turn one
cannot contradict the one an attacker reads on turn five, and the database
banner cannot disagree with the port the same host just advertised. It
must be \emph{safe to serve}: a model that complies too well will invent
credential material with no way to distinguish what it fabricated from
what it leaked, poisoning every downstream analytical product. And it
must be \emph{available}: a generator is a rented service, and ours was.
Operational traces in this study show catalog drift, quota exhaustion,
gateway outages, and safety-alignment refusals. Unverified generation can
forfeit all three properties in the same session.

The alternative---fixing the problem inside the model---fails on
inspection: prompting, refusal training, and constrained decoding shape
probabilities, but none can certify a particular output. Alignment
training produces refusal behavior that shifts with phrasing and
pressure~\cite{ganguli2022,bai2022}; our quota cliff exposed the same
failure under production conditions. And a deceived honeypot is itself an
LLM application with attacker-controlled context, so injection
techniques~\cite{greshake2023} apply with the attacker already in the
room. A probability, however small, is not a promise a security tool can
serve on; the assurance has to live outside the model.

The classical system-side fixes fail differently. Pre-scripting every
answer regresses to the finite automaton of low-interaction deception,
discarding the adaptivity that motivated generation in the first place;
real-time human review of every response cannot keep pace with an
attacker who reads faster than an analyst can vet. What survives is an
asymmetry between the two halves of the problem. Generating a plausible
answer to an arbitrary question is open-ended: its output space cannot
be enumerated, so no amount of testing certifies the next draw.
Checking a draft against a small, closed set of claims is the opposite:
a bounded computation whose verdict is final. A system can exploit this
by splitting the roles, letting the generator propose and a verifier
dispose against a fiction the system itself fully controls.

This paper takes the position that the right question is not
\emph{can an LLM impersonate a host} but \emph{how much intelligence can
a lie safely spend, and where must the spending stop}. Deception systems
build intelligence in rungs -- templates, classifiers, retrieval, free
generation -- and each rung has a trust price that practice pays
implicitly. We make the pricing explicit. Our answer is a verified
deception loop in which
each mechanism exists to buy one of the three properties, placed on a
single axis of certifiability (\S\ref{sec:ladder}). A fiction graph
binds every host fact to one consistent set of triples, and a verifier
checks every draft against that graph before anything reaches the wire:
consistency, enforced at serving time, not asked for at prompt time.
Baited content is drawn only from canary-registered secrets, so that any
disclosure is a measured event rather than a leak: safety by
construction. A deterministic fallback serves human-vetted transcripts
whenever generation or verification fails, so the system degrades to the
last known good answer rather than an unverified improvisation:
availability under provider failure. Beneath these sit a behavioral
profiler and a hybrid retrieval engine that ground a per-attacker
deception plan in curated CTI, giving the lie its targeting, and a
human-in-the-loop gate holds back any content class marked as requiring
approval. The loop turns generation from an act of trust into a measured
pipeline with an audit trail, and the evaluation reads the whole
apparatus as an escalation study: each rung of intelligence is measured
for what it adds over the rung beneath it, including the rungs that fold
back down when a certifiable mechanism reproduces their effect.

We evaluate whether each element of this loop contributes measurably,
using
artifact-backed benchmarks that we release alongside the paper. Four
questions structure the study:
\begin{itemize}
\item \textbf{RQ1:} Does retrieval grounding change CTI analysis, and
separately what is served? A paired ablation across 90 matched scenarios
reports grounding lifts analysis-channel technique detection by $+0.317$
(95\% CI $[+0.162, +0.480]$) and actor attribution by $+0.140$
($[+0.083, +0.197]$), leaving persona realism untouched---a null the
architecture predicts; the served-text comparison is reported separately,
never causally.
\item \textbf{RQ2:} Can an LLM-driven honeypot remain internally
consistent and safe under adversarial interaction? Across 45 targeted
re-queries the fiction graph serves zero contradictions; across a
218-query campaign the verification layer records no unverified response
served; and three authored model families expose bounded
precision/recall gaps.
\item \textbf{RQ3:} Is the resulting deception strong enough to hold a
skilled observer? A three-evaluator consensus study (independently
prompted LLM reviewer personas sharing one judge model) scores the full
session transcripts at $3.86/5$ with exact agreement $0.750$, all eight
scenarios at or above the acceptance bar, and a known judge bias
corrected by canary-registered secrets.
\item \textbf{RQ4:} Does each additional rung of intelligence justify its
trust price? The escalation analysis adds each mechanism to the same loop
and prices its marginal effect: on one 400-session held-out split the
profiler scores $0.8875$ as a whole-session classifier and $0.955$ when
the trajectory replay exposes each session's full command history---the
same corpus under two procedures, not an intervention delta; grounding
contributes the $+0.154$ comparison above; and kind-aware classification corrects 12/13 Cowrie-log assessments and
lifts bait exposure 125\% in five replays yet adds no measurable dwell
and is subsumed by a deterministic floor---a fold-down signal turning the
ladder into a design rule. A second pre-registered contrast on the serving surfaces---the first
between deception surfaces, to our knowledge---repeats the discipline:
both arms saturate the thirty-turn cap (log-rank $p=0.975$) and bait
capture is zero in both; the one separating interval (planted artifacts on adversarial reads,
11/20 versus 0) was promoted to the primary endpoint of a third
pre-registered contrast with a family-disjoint brain and replicated
confirmatorily (6/10 versus 0/10, $p=0.011$).
\end{itemize}

\noindent Our contributions are:
\begin{itemize}
\item \textbf{A verified deception architecture} (\S\ref{sec:design})
coupling RAG-grounded generation to a fiction graph and a serving-time
verification gate with deterministic fallback: every response is verified
against the deception state or drawn from vetted content.
\item \textbf{An archived analysis-channel grounding comparison}
(\S\ref{sec:evaluation}): paired-difference bootstrap intervals over
identical 90-row scenario sets report $+0.154$ overall improvement
($[+0.092, +0.221]$), with 35 TTP detections spanning 21 distinct techniques
versus zero without grounding. No clean served-text causal gain is claimed:
our own scorer audit yields echo-corrected lower bounds of 0.479 technique
and 0.243 actor for the archived campaign.
\item \textbf{A consistency and safety benchmark} (\S\ref{sec:evaluation}):
contradiction-free serving under adversarial re-query, a supported-claim
gate, a 105-case deterministic suite, a cross-family stress test, and an
ablation battery in which removing each mechanism (A1--A4) costs the
property it was added to provide.
\item \textbf{Run-to-run stability}: two archived 218-query reruns
reproduce the composite to $+0.005$ with $0.991$ per-row pass agreement
and exact deterministic-strata reproduction.
\item \textbf{A release package} on commodity hardware
(2~vCPU, ${\sim}4$~GB RAM), in which every headline number resolves to a
persisted result file, published on artifact acceptance.
\item \textbf{An escalation methodology for AI in deception}
(\S\ref{sec:ladder}, \S\ref{sec:jeveval}): mechanisms placed on a
certifiability axis, each rung priced by measured marginal effect under
pre-registered protocols, and the descent direction exercised on live
machinery -- when a certifiable mechanism reproduces an intelligent
one's benefit, the intelligent rung is removed from the ladder; a frozen
candidate that added no value beyond the gate on a prospective 40-case
cohort was retained disabled rather than claimed.
\end{itemize}

We are careful about what this paper does not claim. The deployment is
tailnet-only; the evaluation judge shares a provider with, but is a
different model from, the generator; the only labeling lane in this
release is the disclosed AI panel of \S\ref{sec:labels}; and the
full-system study that exercises the actual serving path ran in
deterministic provider-unavailable mode only. The threat model excludes
host compromise, provider compromise, and operator collusion, and the
closed command surface does not establish open-world robustness. These
limits are quantified where they arise (\S\ref{sec:threat-model},
\S\ref{sec:conclusion}).

Section~\ref{sec:threat-model} fixes the threat model and claim
boundary; \S\ref{sec:design} presents the loop; \S\ref{sec:related}
positions it; \S\ref{sec:evaluation} reports the battery, ablations, and
stability runs; \S\ref{sec:conclusion} concludes.
